% Generated draft: requires mathematical and source review.
\documentclass[11pt]{article}
\usepackage[margin=25mm]{geometry}
\usepackage{fontspec}
\setmainfont{Latin Modern Roman}
\usepackage{amsmath,amssymb,mathtools,unicode-math}
\setmathfont{Latin Modern Math}
\usepackage{enumitem,needspace,xurl,hyperref}
\setlength{\parindent}{0pt}
\setlength{\parskip}{5pt}
\setlength{\emergencystretch}{3em}
\newcommand{\N}{\mathbb{N}}
\newcommand{\Z}{\mathbb{Z}}
\newcommand{\Q}{\mathbb{Q}}
\newcommand{\R}{\mathbb{R}}
\newcommand{\C}{\mathbb{C}}
\newcommand{\F}{\mathbb{F}}
\newcommand{\A}{\mathbb{A}}
\newcommand{\PP}{\mathbb{P}}
\newcommand{\E}{\mathbb{E}}
\newcommand{\eps}{\varepsilon}
\newcommand{\dd}{\,\mathrm{d}}
\begin{document}
\section*{UnsolvedMath ID: 30002555}
\textbf{Source problem code:} catalogue entry
\subsection{Definitions and mathematical statement}
Let \(C\) be the Cantor set and let \(\operatorname{GL}^{+}(2,\R)=\{A\in M_2(\R):\det A>0\}\). A translation surface is understood here as a connected, oriented surface without boundary, equipped away from its singularities with Euclidean coordinate charts whose transition maps are translations. Its translation structure induces a flat metric. The usual tameness convention is that points of the metric completion have local models given by regular Euclidean disks, finite-angle Euclidean cones of angle \(2\pi k\), \(k\geq 1\), or infinite-angle cones. The supplied record does not specify its precise convention for tameness; this local-model formulation requires confirmation against the source.

An orientation-preserving affine automorphism of \(S\) is a self-homeomorphism that, in translation coordinates, has the form \(z\mapsto Az+b\), with a constant derivative \(A\in\operatorname{GL}^{+}(2,\R)\). Its Veech group is
\[
\operatorname{Veech}(S)=\{D f:f\in\operatorname{Aff}^{+}(S)\}\leq\operatorname{GL}^{+}(2,\R).
\]
The realization requested below is equality with the prescribed matrix subgroup \(G\), rather than merely an abstract group isomorphism.

Write \(\operatorname{Ends}(S)\) for the topological end space, defined by compatible choices of components escaping compact subsets of \(S\), with its usual end topology. An end is planar if it has a neighborhood of genus zero. Thus having no planar ends means that every end is accumulated by genus. Infinite genus means that the genera of compact subsurfaces of \(S\) are unbounded.

\textbf{Supplied realization question.} For every closed subspace \(X\subseteq C\) and every countable subgroup \(G\leq\operatorname{GL}^{+}(2,\R)\) without contracting elements, does there exist a tame translation surface \(S\) of infinite genus, without planar ends, such that
\[
\operatorname{Veech}(S)=G
\qquad\text{and}\qquad
\operatorname{Ends}(S)\cong X,
\]
where \(\cong\) denotes homeomorphism of topological spaces?

\textbf{Unresolved qualifications.} The record does not define ``contracting.'' A conventional metric interpretation is that \(A\) is contracting when
\[
\|A\|_{\mathrm{op}}:=\sup_{\|v\|_2=1}\|Av\|_2<1,
\]
so that the hypothesis becomes \(\|A\|_{\mathrm{op}}\geq1\) for every \(A\in G\). This interpretation is provisional and must be checked against the cited report. Also, the stated condition that \(X\) be closed permits \(X=\varnothing\). Under the usual surface conventions, a connected surface of infinite genus is noncompact and has a nonempty end space. Consequently, the empty case is incompatible with the requested realization. The supplied record does not state whether \(X\neq\varnothing\) was intended; that restriction has not been silently inserted.
\subsection{Short English statement}
The question asks whether every countable matrix group satisfying the source's noncontraction condition occurs as the Veech group of a tame surface of infinite genus whose end space is prescribed and whose ends are all accumulated by genus. The empty end-space case and the precise contraction and tameness conventions require clarification. The supplied status is partially solved; this draft does not independently verify that assessment.
\subsection{Sources}
\begin{itemize}
\item Supplied local UnsolvedMath registry record, permanent ID \(30002555\), title \emph{Veech Groups with Prescribed End Spaces}, source problem code \texttt{catalogue entry}, in the collection \emph{Oberwolfach Reports Open Problems}. The proposed year is recorded as 2014; no individual proposer is supplied.
\item Supplied primary report citation: \emph{Flat Surfaces and Dynamics on Moduli Space} (2014), Oberwolfach Reports, pp.\ 869--941, DOI \texttt{10.4171/owr/2014/15}. \url{https://doi.org/10.4171/owr/2014/15}. The background describes this report as documenting workshop discussions and open problems.
\item Supplied literature reference: Przytycki--Schmithusen--Valdez, \emph{Annales de l'Institut Fourier}. \url{https://doi.org/10.5802/aif.2625}.
\item Supplied literature reference: Ramirez Maluendas--Valdez, arXiv. \url{https://arxiv.org/abs/1603.00503}.
\item Supplied literature assessment, checked on \texttt{2026-08-21}: partially solved. The supplied assessment says that realization is known for the one-ended Loch Ness monster and the blooming Cantor tree, while realization for an arbitrary prescribed closed end space remains open. This is supplied dated curation derived from the OWR material, not an independently verified literature assessment; time-sensitive status and the cited claims require independent verification.
\end{itemize}
\end{document}
